\scanpage{110}
\mspage{97}
% archive sheet number: 109

d'où
\[
 u_{!}(TC)^{\sim}\subset TC'
\]
puisque $TC'$ est Q-clos à gauche, \deleted{\unsure{(ceci}} \inserted{(ou parce
que ça a été prouvé dans (5.8.14))}, \deleted{résulte d'autre part} on a \unsure{déjà} prouvé l'inclusion en sens inverse (5.8.12 bis),
\unsure{ess.} équivalente à 5.8.4 (et qui représentait
la principale difficulté).

\uline{Remarque} 5.9. Si on ne suppose pas que la
\deleted{soit} structure sur $M$ soit close, alors
la démonstration donnée de 5.8. montre
qu'on trouve \deleted{\unsure{quand}} sur $M'$ une
structure de Quillen spéciale, en définissant
cette fois $C'$ comme dans 4.6. Mais je
doute qu'on aura \inserted{jamais} à travailler avec des
triples \unsure{de} Quillen spéciaux qui ne soient
clos. De toute façon, si $W$ est
stable par petites $\varinjlim$ filtrantes\textsuperscript{$*$},
% slanted
\margin{ce qui \unsure{implique} la stabilité par facteurs directs, et \unsure{a fortiori|à fortiori} \unsure{aussi} la $=$ propriété pour \deleted{\unsure{$u^{*-1}(W)$}} $W'\overset{\mathrm{déf}}=(u^{*})^{-1}(W)$.}
\deleted{\unsure{donc}} \inserted{(et je} doute qu'il y \unsure{ait} \deleted{un triple de Quillen spécial} \inserted{des triples de Quillen spéciaux}
en dehors de ce cas), alors si
$(W,C,F)$ est un triple de Quillen
spécial, il en est de $=$ de $(W,\widetilde C=C^{\unsure{\natural|4}},F)$,
et je présume qu'on \deleted{\unsure{très}} aurait intérêt
à remplacer \unsure{le|ce} triple initial par celui-là.

